\chapter*{Preface}
\addcontentsline{toc}{chapter}{Preface}

Computational environmental science increasingly depends on the orchestration
of heterogeneous resources: numerical models compiled for HPC clusters, data
preprocessing and post-processing scripts, container images that freeze a
fragile scientific software stack, cloud instances rented for a single
forecast cycle, and --- more recently --- AI services consumed through web
APIs for summarization, anomaly detection, or decision support. The challenge
this book addresses is not how to write any one of these components, but how
to make their \emph{composition} explicit, reproducible, portable, and
auditable.

\emph{DAGonStar} answers this challenge with a deliberately small abstraction:
a workflow is a directed acyclic graph (DAG) of tasks, and a task is a Python
object that executes a command --- local, remote, containerized,
scheduler-submitted, an HTTP request, an LLM completion, or an importable
Python function --- in a particular environment. This minimalism is the
engine's central design virtue. It keeps workflow authoring close to ordinary
Python, while still providing the scientific guarantees that ad hoc shell
scripts cannot: explicit dependency graphs, automatic data staging, cycle
detection, checkpointing, and a uniform mechanism --- the \texttt{workflow://}
schema --- for expressing that one task consumes a file produced by another.

\paragraph{How this book is organized.}
The book accompanies a course of 48 lectures plus one appendix lecture. Each
lecture has a slide deck of 30 slides (\texttt{en/lectures/}), speaker notes
(\texttt{en/notes/}) and one or more runnable examples
(\texttt{examples/NN-topic/}). Theory and hands-on work are not separated:
every lecture introduces a concept and runs it immediately. The chapters
group the lectures into four parts, and \Cref{tab:course-map} shows where
each lecture is treated.

\begin{table}[h]
\centering
\small
\begin{tabular}{@{}lp{0.47\linewidth}p{0.30\linewidth}@{}}
\toprule
\textbf{Lectures} & \textbf{Topics} & \textbf{Chapters} \\
\midrule
01--05 & Workflows, script vs DAG, architecture, environment, configuration &
  \ref{ch:introduction}--\ref{ch:getting-started}, \ref{ch:configuration} \\
06--10 & First task, DAGs, dataflow, cycles, scratch folders &
  \ref{ch:lessons-01-02}, \ref{ch:schema}, \ref{ch:lessons-03-04} \\
11--12 & Inspection, station pipeline & \ref{ch:course-pipeline} \\
13, 15 & Staging, checkpoints & \ref{ch:lessons-05-07} \\
14, 16, 19--21, 23--24 & Garbage collection, failures, bindings, sweeps,
  ensembles, dry runs & \ref{ch:resilient} \\
17, 22 & Asynchronous runs, meta-workflows & \ref{ch:meta-async} \\
18 & Native tasks & \ref{ch:native-tasks} \\
25--28 & Web and LLM tasks & \ref{ch:web-tasks}, \ref{ch:llm-tasks} \\
29--36 & FaaS, IoT, agentic tasks, continuum & \ref{ch:faas-iot-agentic} \\
37, 41 & Docker, Slurm & \ref{ch:backends} \\
38--40, 42--44 & Apptainer, orchestrators, remote and cloud, BeeOND,
  pre-staging, movers & \ref{ch:hpc-storage} \\
45--48 & FAIR, CWL, operations, capstone & \ref{ch:fair-cwl} \\
A1 & DAGonFS & Appendix~\ref{app:dagonfs} \\
A2 & Ad-hoc volumes and streams & Appendix~\ref{app:volumes} \\
A3 & Meta-tasks & Appendix~\ref{app:meta-tasks} \\
A4 & DAGonWeb & Appendix~\ref{app:dagonweb} \\
\bottomrule
\end{tabular}
\caption{Course map: lectures and the chapters that treat them.}
\label{tab:course-map}
\end{table}

Chapters on external software (\Cref{ch:external-software}), the
environmental case study (\Cref{ch:case-study}), the examples catalog
(\Cref{ch:examples-catalog}) and the synthesis (\Cref{ch:synthesis}) connect
the lectures to research practice. Appendices~A--D collect reference
material: the public API, developer practices, the configuration model and
checkpoint semantics. Appendix~E presents DAGonFS, Appendix~F ad-hoc volumes
and streams, Appendix~G meta-tasks and Appendix~H DAGonWeb.

\paragraph{Intended audience.}
This book is written for graduate students, researchers, and research software
engineers in computational and environmental science who need to compose
heterogeneous computational pipelines reproducibly. Working knowledge of
Python and of basic Unix/HPC concepts (shells, SSH, job schedulers) is
assumed; no prior exposure to workflow engines is required.

\paragraph{A note on fidelity to the software.}
Every course example is included in this book with \texttt{lstinputlisting}
from the course repository, and so is its recorded output. The examples were
written for and checked against DAGonStar~3.0.0
(\url{https://github.com/DagOnStar/dagonstar}). The script
\texttt{tools/run\_examples.py} runs them all again and reports any output
that differs from the recorded one. Where the software has a limitation,
such as name-based checkpoint reuse (\Cref{ch:resilient}), the book shows it
rather than smoothing it away.
